% Nouveau EGA V (cote 25): archive cote 25, site page 127 (PDF page 128).
% typescript. Draft compared: none.
% AI-transcribed by Claude Opus 5.5 (claude-opus-5-5, adaptive thinking, effort high), Message Batches API, batch ega25b (msgbatch_01A7n5YuXubUPNK1zsepi8RG), 2026-11-13; not yet checked by eye.
\scanpage{25/127}
\mspage{30}
brique de k(s), il existe un cycle premier T associé à $X_k$ tel que
$k \to H^0(T,\underline{O}_T)$ soit un isomorphisme (condition automatiquement satisfaite
si g est propre et surject\uline{i}f). Alors le morphisme D: $\check{P} \longrightarrow \mathrm{Div}_{X/S}$ est
un monomorphisme.

b) \xout{Supposonsxxf} Considérons l'application $u \rightsquigarrow D\circ u$ de
$\mathrm{Aut}_S(\check{P}\ )$ dans l'ensemble des familles de diviseurs sur X/S paramétrées
par P . Alors si \unsure{g} est surjectif, l'application précédente est injective,
en particulier D = D u implique u = $\mathrm{id}_{\check{P}}$\inserted{; plus généralement, le morphisme de foncteurs $\underline{\mathrm{Aut}}_S(\check{P}) \longrightarrow \underline{\mathrm{Syst.\ lin.\ div}}_{X/S,\check{P}}$ est un monomorphisme}

On notera que sous leshypothèses de a), b) est uneconséquence
triviale de a); d'autre part, b) est valable sous d eshypothèses moins
restrictives que a). On fera attention que a) devient faux si on abandon
ne l'hypothèse restrictive qui y a \xout{xéé} été faite: prendre par exemple
S=Spec(k), X un ouvert de $P^1_k$ ne contenant pas les deux points distincts
a et b de $P^1_k(k)$ , D la famille de diviseurs définie par l'immersion
canonique de X dans $P^1_k$ . Alors les deux \xout{morphismes} points a et b
définissent les mêmes diviseurs de X (savoir le diviseur nul), sans être
eux-mêmes identiques.

Supposons d'abord que S soit le spectre d'un corps \inserted{k, qu'on peut évidemment (par ``descente'') supposer} alg clos. \deleted{\illegible{1}}
Soit T comme dans a)\inserted{, \unsure{qu'on munit|qui muni} de la structure réduite induite,} on a alors un morphisme ``diviseur induit''
$\mathrm{Div}_{X/k} \longrightarrow \mathrm{Div}_{T/k}$ , et il suffit de prouver que le composé $\check{P} \longrightarrow \mathrm{Div}_{T/k}$
est un monomorphisme. Comme ce dernier est encore un syst lin de div ,
on est ramené au cas où X=T , donc au cas où $H^0(X,\underline{O}_X) \xleftarrow{\sim} k$ . Alors pour
tout S sur k, on a $g_{S*}(\underline{O}_{X_S}) \xleftarrow{\sim} \underline{O}_S$ , donc (si $\underline{L}$ sur X et u: $\underline{E} \to$ \deleted{\illegible{1}}$\hand{g_*}(\underline{L})$
sont comme dans §.1.) si deux sections $\phi$, $\psi$ de $\underline{E}_S$ partout non nulles
sont telles que $u(\phi)$ et $u(\psi)$ sont des sections de $\underline{L}_S$ sur $X_S$ ayant même
diviseur, donc \xout{xxxxxf} déduites l'une de l'autre par multiplication par
une section \inserted[typed]{inversible} de $\underline{O}_{X_S}$, il s'ensuit que $\psi$ est déduit de $\phi$ par multiplication

\margin{127}
\margin{\unsure{Mais|Non} automatiquement vérifiée si les $X_s$ sont \deleted{\illegible{1}} intègres.}
